\subsection{分式模的性质}

在本节中，A 表示一个环，S 表示 A 的一个乘性子集。

设 \((\mathbf{M}_{\alpha}, \phi_{\beta\alpha})\) 是 A-模的一个正向系；则 \((\mathrm{S}^{-1}\mathbf{M}_{\alpha}, \mathrm{S}^{-1}\phi_{\beta\alpha})\) 是 \(S^{-1}A\)-模的一个正向系，而取正向极限与张量积交换这一事实（《代数学》第 II 章，§ 6，第 3 节，命题 7）使我们可以定义一个典范同构

\[
\lim _ {\longrightarrow} \left(\mathrm{S} ^ {- 1} \mathrm{M} _ {\alpha}\right)\rightarrow \mathrm{S} ^ {- 1} \lim _ {\longrightarrow} \mathrm{M} _ {\alpha}.
\]

类似地，取直和与张量积交换这一事实（《代数学》第 II 章，§ 3，第 7 节，命题 7）使我们可以对 A-模的每个族 \((\mathbf{M}_t)_{i\in I}\) 定义一个典范同构

\[
\bigoplus_ {i \in I} S ^ {- 1} M _ {i} \rightarrow S ^ {- 1} \bigoplus_ {i \in I} M _ {i}.
\]

最后我们指出，若 A-模 M 是子模族 \((\mathbf{N}_{\iota})_{\iota\in I}\) 的和，则 \(S^{-1}M\) 是诸子 \(S^{-1}A\)-模之和，这些子模由 \(i_{M}^{S}(N_{\iota})\) 生成。由此得出，若 M 是有限生成的 A-模（相应地，有限表示的 A-模），则 \(S^{-1}M = S^{-1}A \otimes_{A} M\) 是有限生成的 \(S^{-1}A\)-模（相应地，有限表示的 \(S^{-1}A\)-模）。

定理 1. 环 \(\mathbf{S}^{-1}\mathbf{A}\) 是一个平坦的 A-模（第 I 章，§ 2，第 1 节，定义 2）。

若 \(u: \mathbf{M}' \to \mathbf{M}\) 是 A-模的单同态，则须证明 \(\mathbf{S}^{-1}u: \mathbf{S}^{-1}\mathbf{M}' \to \mathbf{S}^{-1}\mathbf{M}\) 是单射。现在，若 \(m'/s\)（\(m' \in \mathbf{M}'\)、\(s \in \mathbf{S}\)）使得 \(u(m')/s = 0\)，这就意味着存在 \(s' \in \mathbf{S}\) 使得 \(s'u(m') = 0\)（第 2 节，命题 4），亦即 \(u(s'm') = 0\)；由于 \(u\) 是单射，得出 \(s'm' = 0\)，从而 \(m'/s = 0\)。

\(\mathbf{S}^{-1}\mathbf{A}\) 是平坦 A-模这一事实使我们可以对它应用第 I 章 § 2 的结果。特别地：

\begin{enumerate}
\def\labelenumi{(\arabic{enumi})}
\tightlist
\item
  若 M 是 A-模而 N 是 M 的子模，则 \(S^{-1}N\) 典范地等同于 \(S^{-1}M\) 的由 \(i_{\mathrm{M}}^{\mathrm{S}}(\mathrm{N})\) 生成的子模（第 I 章，§ 2，第 3 节，注 2）；在这一等同下，\(S^{-1}(\mathrm{M}/\mathrm{N})\) 与 \((\mathrm{S}^{-1}\mathrm{M})/(\mathrm{S}^{-1}\mathrm{N})\) 等同，并且，若 P 是 M 的第二个子模，则
\end{enumerate}

\[
\mathrm{S} ^ {- 1} (\mathrm{N} + \mathrm{P}) = \mathrm{S} ^ {- 1} \mathrm{N} + \mathrm{S} ^ {- 1} \mathrm{P}, \quad \mathrm{S} ^ {- 1} (\mathrm{N} \cap \mathrm{P}) = \mathrm{S} ^ {- 1} \mathrm{N} \cap \mathrm{S} ^ {- 1} \mathrm{P}
\]

（第 I 章，§ 2，第 6 节，命题 2）。

\begin{enumerate}
\def\labelenumi{(\arabic{enumi})}
\setcounter{enumi}{1}
\tightlist
\item
  若 M 是有限生成的 A-模，则
\end{enumerate}

\[
\mathbf {S} ^ {- 1} \operatorname{Ann} (\mathbf {M}) = \operatorname{Ann} (\mathbf {S} ^ {- 1} \mathbf {M})\tag{9}
\]

（第 I 章，§ 2，第 10 节，命题 12 的推论 2）。

命题 10. 设 M 是 A-模。对 S \(^{-1}\) A-模 S \(^{-1}\) M 的每个子模 N′，用 \(\phi(N')\) 表示 N′ 在 \(i_{M}^{8}\) 下的原像。则：

\begin{enumerate}
\def\labelenumi{(\roman{enumi})}
\item
  \(\mathbf{S}\phi(\mathbf{N}^{\prime}) = \mathbf{N}^{\prime}.\)
\item
  对任一子模 \(\mathbf{N}\)（\(\mathbf{M}\) 的），子模 \(\phi (\mathrm{S}^{-1}\mathrm{N})\)（\(\mathbf{M}\) 中的）由这样的 \(m\in \mathbf{M}\) 组成：存在 \(s\in \mathbf{S}\) 使得 \(sm\in \mathbf{N}\)。
\item
  \(\phi\) 是子 \(\mathbf{S}^{-1}\mathbf{A}\)-模（\(\mathbf{S}^{-1}\mathbf{M}\) 的）的集合到满足下列条件的子模 \(\mathbf{Q}\)（\(\mathbf{M}\) 的）组成的集合上的一个同构（关于由包含关系定义的序）：

(MS) 若 \(sm \in \mathbb{Q}\)，其中 \(s \in \mathbb{S}\)、\(m \in \mathbb{M}\)，则 \(m \in \mathbb{Q}\)。

\end{enumerate}
显然 \(\mathrm{S}^{-1}\phi(\mathrm{N}')\subset\mathrm{N}'\)；反之，若 \(n' = m/s \in N'\)，则 \(m/1 \in N'\)，从而 \(m \in \phi(\mathrm{N}')\)，于是 \(n' \in \mathrm{S}^{-1}(\phi(\mathrm{N'}))\)；由此得 (i)。为使元素 \(m \in M\) 满足 \(m \in \phi(\mathrm{S}^{-1}\mathrm{N})\)，必须且只需 \(m/1 \in S^{-1}N\)，即存在 \(s \in S\) 与 \(n \in N\) 使得 \(m/1 = n/s\)；这就意味着存在 \(s' \in S\) 使得 \(s'sm = s'n \in N\)，由此得 (ii)。最后，关系 \(sm \in \phi(\mathrm{N}')\) 按定义等价于 \(sm/1 \in N'\)，而由于 s/1 在 \(S^{-1}A\) 中可逆，这就蕴含 \(m/1 \in N'\)，亦即 \(m \in \phi(\mathrm{N}')\)，因而 \(\phi(\mathrm{N}')\) 满足条件 (MS)；另一方面，由 (ii) 可知，若 N 满足 (MS)，则 \(\phi(\mathrm{S}^{-1}\mathrm{N}) = N\)，这就完成了 (iii) 的证明。

子模 \(\phi(\mathrm{S}^{-1}\mathrm{N})\) 称为 N 在 M 中关于 S 的饱和，而满足条件 (MS)（从而等于其饱和）的子模称为关于 S 是饱和的。子模 \(\phi(\mathrm{S}^{-1}\mathrm{N})\) 是下列复合同态的核

\[
\mathrm{M} \xrightarrow {h} \mathrm{M} / \mathrm{N} \xrightarrow {i _ {\mathrm{M} / \mathrm{N}} ^ {\mathrm{S}}} \mathrm{S} ^ {- 1} \mathrm{M} / \mathrm{S} ^ {- 1} \mathrm{N}
\]

其中 h 是典范同态，这由图表

\[
\begin{array}{c c c} \mathrm{M} & \stackrel {{h}} {{\longrightarrow}} & \mathrm{M/N} \\ i _ {\mathrm{M}} ^ {\mathrm{S}} \Big \downarrow & & \Big \downarrow i _ {\mathrm{M/N}} ^ {\mathrm{S}} \\ \mathrm{S} ^ {- 1} \mathrm{M} & \xrightarrow [ \mathrm{S} ^ {- 1 h} ]{} \mathrm{S} ^ {- 1} \mathrm{M/S} ^ {- 1} \mathrm{N} \end{array}
\]

若 S 是 A 中素理想 p 的余集，\(\phi(S^{-1}N)\) 也称为 N 在 M 中关于 p 的饱和。

推论 1. 设 \(N_{1}\)、\(N_{2}\) 是 A-模 M 的两个子模。为使 \(S^{-1}N_{1} \subset S^{-1}N_{2}\)，必须且只需 \(N_{1}\) 关于 S 的饱和包含于 \(N_{2}\) 的饱和。

推论 2. 若 M 是 Noether（相应地，Artin）A-模，则 S \(^{-1}\) M 是 Noether（相应地，Artin）S \(^{-1}\) A-模。特别地，若环 A 是 Noether 的（相应地，Artin 的），则环 S \(^{-1}\) A 亦然。

